\subsection*{Sub-Level Parallelism}

Targets like  allow a single parallelism level to be divided into up to three sub-levels. In Croqtile, \texttt{ituple} and \texttt{mdspan} are used for this syntax:

\begin{lstlisting}[language=croqtile,caption={Sub-Level Parallelism}]
parallel {px, py, pz} by [1, 1, 2] {
  // higher synchronization cost
  // ...
  parallel t_index = {qx, qy} by [3, 4] {
    // lower synchronization cost
    // ...
  }
  // ...
}
\end{lstlisting}

In this example, two levels of parallelism are defined. The first level (\emph{level-0}) is subdivided into three sub-levels, annotated by the \texttt{\{px, py, pz\}} \texttt{ituple}. The second level (\emph{level-1}) is subdivided into two sub-levels, annotated by the \texttt{\{qx, qy\}} \texttt{ituple}, or \texttt{t\_index} as a whole. For each sub-level, the parallelization count corresponds to the respective value inside the \texttt{mdspan}.

The first element inside the \texttt{ituple} represents the \textbf{least significant parallel variable}. Croqtile adopts this ordering to align with  developers' conventions, rather than for internal consistency.
